\section{Idea de negocio}

\subsection{Servicio ofrecido}

GymHub ofrece acceso a una plataforma web para organizar la operación de un
gimnasio. El servicio permite registrar miembros, gestionar planes de membresía,
controlar pagos y asistencia, asignar planes de entrenamiento, registrar sesiones
y consultar indicadores básicos. La aplicación no constituye un dispositivo
médico ni un sistema contable certificado.

La versión inicial se concentra en funciones comprobables del prototipo. Las
integraciones con pasarelas de pago, acceso físico, almacenamiento persistente de
archivos y automatizaciones de producción se consideran fases posteriores.

\subsection{Cliente y personas usuarias}

El cliente principal es el dueño o administrador de un gimnasio pequeño o
mediano, porque asume el costo y define la adopción. Como segmento complementario
se consideran entrenadores personales que administran una cartera propia. Los
miembros del gimnasio son beneficiarios y usuarios finales, pero no necesariamente
son compradores directos.

El mercado inicial se delimita a Pérez Zeledón y zonas cercanas para facilitar
demostraciones, acompañamiento y retroalimentación. El número de establecimientos
que cumplen el perfil es \pendiente{} de determinar mediante investigación de
campo y fuentes locales verificables.

\subsection{Forma de comercialización}

La adquisición se plantea mediante contacto directo, demostración del prototipo,
prueba piloto con alcance limitado y posterior suscripción. Los canales de apoyo
serán una página informativa, WhatsApp Business, correo y contenido demostrativo
en redes sociales. No se proyectará publicidad pagada antes de validar el mensaje
y la disposición de pago.

\subsection{Necesidad atendida}

La propuesta busca reducir la dispersión de información y facilitar el control
diario de clientes, cobros, vigencias, asistencia y programas de entrenamiento.
El valor no consiste únicamente en digitalizar un cuaderno, sino en relacionar
datos para mostrar acciones: pagos vencidos, miembros inactivos, entrenamiento
del día y progreso registrado.

\subsection{Impacto esperado}

Se espera mejorar la organización del gimnasio, la trazabilidad de operaciones y
la comunicación con miembros. El impacto ambiental potencial corresponde a la
reducción de registros impresos. Ambos efectos deberán medirse durante el piloto
con indicadores como tiempo de registro, incidencias administrativas, formularios
eliminados y frecuencia de uso.

\section{Modelo de negocio}

El modelo consiste en software como servicio por suscripción mensual. GymHub
mantiene la aplicación y el gimnasio adquiere derecho de uso, actualizaciones y
un nivel de soporte según el plan. La estructura sigue cuatro dimensiones:

\begin{itemize}
  \item \textbf{Cliente}: responsables de gimnasios y entrenadores con necesidad
  de centralizar información.
  \item \textbf{Oferta}: administración y seguimiento deportivo en una sola
  plataforma responsive.
  \item \textbf{Infraestructura}: equipo de desarrollo, aplicación, base de datos,
  servicios de nube, soporte y procesos de seguridad.
  \item \textbf{Viabilidad económica}: suscripciones recurrentes que deben cubrir
  infraestructura, soporte, mantenimiento e inversión de trabajo.
\end{itemize}

El Canvas permite describir cómo se crea, entrega y captura valor
\parencite{osterwalder2010}. En GymHub, sus nueve bloques se complementan con
evidencia técnica y con una validación comercial pendiente, en lugar de tratar el
lienzo como prueba suficiente de viabilidad.

\begin{table}[H]
\caption{Resumen del modelo de negocio de GymHub}
\begin{tabularx}{\textwidth}{>{\bfseries}p{3.5cm}X}
\toprule
Elemento & Definición inicial \\
\midrule
Servicio & Plataforma web de gestión administrativa y seguimiento deportivo. \\
Comprador & Dueño o administrador de gimnasio pequeño o mediano. \\
Usuarios & Administradores, entrenadores, personal autorizado y miembros. \\
Ingreso & Suscripción mensual y, en fases posteriores, configuración o soporte especializado. \\
Entrega & Acceso web responsive con acompañamiento inicial. \\
Mercado inicial & Pérez Zeledón y zonas cercanas; tamaño pendiente de validar. \\
Madurez & MVP funcional en demostración, todavía no SaaS comercial maduro. \\
\bottomrule
\end{tabularx}
\caption*{\fuentepropia}
\end{table}
